\documentclass[12pt,a4paper]{article}

% ── Packages ────────────────────────────────────────────────────────────────
\usepackage[margin=2.5cm]{geometry}
\usepackage{amsmath, amssymb}
\usepackage{booktabs}
\usepackage{xcolor}
\usepackage{listings}
\usepackage{hyperref}
\usepackage{parskip}
\usepackage{microtype}
\usepackage{titlesec}
\usepackage{enumitem}
\usepackage{graphicx}
\usepackage{float}

% ── Colour scheme ────────────────────────────────────────────────────────────
\definecolor{fifablue}{RGB}{0,55,140}
\definecolor{codegreen}{RGB}{0,128,0}
\definecolor{codegrey}{RGB}{128,128,128}
\definecolor{codebg}{RGB}{248,248,248}
\definecolor{codered}{RGB}{180,0,0}

% ── Section formatting ───────────────────────────────────────────────────────
\titleformat{\section}{\large\bfseries\color{fifablue}}{}{0em}{}[\titlerule]
\titleformat{\subsection}{\normalsize\bfseries\color{fifablue!80!black}}{}{0em}{}

% ── Code listings ────────────────────────────────────────────────────────────
\lstset{
  language=Python,
  backgroundcolor=\color{codebg},
  basicstyle=\ttfamily\small,
  keywordstyle=\bfseries\color{fifablue},
  commentstyle=\itshape\color{codegrey},
  stringstyle=\color{codered},
  showstringspaces=false,
  breaklines=true,
  frame=single,
  rulecolor=\color{codegrey!40},
  numbers=left,
  numberstyle=\tiny\color{codegrey},
  numbersep=8pt,
  xleftmargin=1.5em,
  framexleftmargin=1.5em,
}

% ── Hyperlink setup ──────────────────────────────────────────────────────────
\hypersetup{
  colorlinks=true,
  linkcolor=fifablue,
  urlcolor=fifablue,
  citecolor=fifablue,
}

% ── Math convenience macros ──────────────────────────────────────────────────
\newcommand{\Poi}{\operatorname{Poisson}}
\newcommand{\E}{\mathbb{E}}
\newcommand{\Prob}{\mathbb{P}}

% ────────────────────────────────────────────────────────────────────────────
\begin{document}

\begin{titlepage}
  \centering
  \vspace*{2cm}
  {\color{fifablue}\rule{\linewidth}{1pt}}\\[0.4em]
  {\huge\bfseries\color{fifablue} World Cup Predictor}\\[0.4em]
  {\large A Technical Explanation of the Model, Mathematics, and Code}\\[0.2em]
  {\color{fifablue}\rule{\linewidth}{1pt}}\\[1.5em]
  {\large Ben Jordan \& Jonathan Glanfield}\\[0.5em]
  {\normalsize \url{https://github.com/JonathanG04157/worldcup}}\\[0.5em]
  {\normalsize 2026 FIFA World Cup --- USA / Canada / Mexico}\\[2cm]
  {\small\itshape This document explains the full pipeline: how team
  strengths are derived from FIFA rankings, how individual matches are
  simulated using a Poisson model, how the tournament bracket is
  constructed and resolved, and how the Monte~Carlo runner aggregates
  100\,000 simulations into win probabilities.}
  \vfill
  {\small \today}
\end{titlepage}

\tableofcontents
\newpage

% ────────────────────────────────────────────────────────────────────────────
\section{Overview}

The predictor is a \emph{Monte~Carlo simulation} of the 2026 FIFA World
Cup.  Each simulation plays out the entire tournament exactly once,
from the group stage through to the final, by sampling random scorelines
from a statistical model.  After running many thousands of simulations
the code counts how often each team wins, reaches the semi-finals, and so
on, thereby estimating win probabilities.

The project consists of four Python modules and two data files:

\begin{center}
\begin{tabular}{lp{9cm}}
\toprule
\textbf{File} & \textbf{Purpose} \\
\midrule
\texttt{AttackDefenceStrength.py} & Converts FIFA ranking points into
attack and defence strength scores. \\
\texttt{MatchSim.py} & Simulates a single match using the Poisson
model; handles extra time and penalties. \\
\texttt{Tournament.py} & Defines the tournament structure---groups,
fixtures, standings, knockout bracket. \\
\texttt{TourSim.py} & Monte~Carlo runner: simulates the tournament
$N$ times and writes output CSVs. \\
\midrule
\texttt{team\_rankings.csv} & FIFA ranking points per team. \\
\texttt{team\_stats.csv} & Derived attack/defence values (written by
\texttt{AttackDefenceStrength.py}). \\
\bottomrule
\end{tabular}
\end{center}

% ────────────────────────────────────────────────────────────────────────────
\section{Step 1 --- Team Strength from FIFA Rankings}
\label{sec:strength}

\subsection{The FIFA Ranking Points System}

FIFA assigns each national team a \emph{points total} based on their
results over the previous four years, with recent matches weighted more
heavily.  The current range across the 48 World Cup participants runs
from roughly 1\,275 (New Zealand) to 1\,877 (Argentina).

\subsection{The Strength Mapping}

Rather than using raw points directly, the code maps each team's points
to a dimensionless \emph{strength} value centred on 1 via a linear
transformation:

\begin{equation}
  \text{strength}(r) = 1 + \frac{r - 1500}{1000}
  \label{eq:strength}
\end{equation}

where $r$ is the team's FIFA points.  A team rated at exactly 1\,500
points receives strength $= 1.0$ (exactly average).  Argentina at
$r \approx 1877$ receives strength $\approx 1.38$, while New Zealand at
$r \approx 1276$ receives strength $\approx 0.78$.

\subsection{Attack and Defence}

Each team receives two derived attributes:

\begin{align}
  a_i &= \text{strength}(r_i) \label{eq:attack}\\
  d_i &= 2 - \text{strength}(r_i) \label{eq:defence}
\end{align}

The constraint $a_i + d_i = 2$ is a deliberate design choice: better
teams have a higher attack score \emph{and} a lower defence score, where
a \emph{lower} defence score means \emph{fewer} goals conceded (see
Section~\ref{sec:goals}).  The midpoint of both scales is 1.0.

\subsection{Code: \texttt{AttackDefenceStrength.py}}

\begin{lstlisting}
import pandas as pd

def rating_to_strength(rating, base=1.0):
    return base + (rating - 1500) / 1000

df = pd.read_csv('team_rankings.csv')

fifa_rankings = dict(zip(df['Team'], df['Points']))
team_stats = {
    team: {
        'attack':  rating_to_strength(rating),
        'defence': 2.0 - rating_to_strength(rating)
    }
    for team, rating in fifa_rankings.items()
}

df['attack']  = df['Points'].apply(rating_to_strength)
df['defence'] = df['Points'].apply(lambda r: 2.0 - rating_to_strength(r))
df.to_csv('team_stats.csv', index=False)
\end{lstlisting}

The module both builds an in-memory dictionary \texttt{team\_stats} and
persists the values to \texttt{team\_stats.csv} for inspection.

% ────────────────────────────────────────────────────────────────────────────
\section{Step 2 --- The Poisson Match Model}
\label{sec:poisson}

\subsection{Why the Poisson Distribution?}

Football goals are \emph{rare, independent events} occurring at a roughly
constant rate over 90 minutes --- exactly the conditions under which the
Poisson distribution is the natural model.  If a team is expected to
score $\lambda$ goals, the probability of them scoring exactly $k$ goals
is:

\begin{equation}
  \Prob(X = k) = \frac{\lambda^k e^{-\lambda}}{k!}, \quad k = 0, 1, 2, \ldots
  \label{eq:poisson}
\end{equation}

Crucially, the home and away goal tallies are drawn \emph{independently},
so a full scoreline $(g_H, g_A)$ has probability

\begin{equation}
  \Prob(g_H, g_A) = \frac{\lambda_H^{g_H} e^{-\lambda_H}}{g_H!}
                   \cdot
                   \frac{\lambda_A^{g_A} e^{-\lambda_A}}{g_A!}.
\end{equation}

\subsection{Computing Expected Goals}
\label{sec:goals}

For a match between a home team $H$ and an away team $A$, the expected
goals for each side are:

\begin{align}
  \lambda_H &= \lambda_0 \cdot a_H \cdot d_A \label{eq:lambdaH}\\
  \lambda_A &= \lambda_0 \cdot a_A \cdot d_H \label{eq:lambdaA}
\end{align}

where $\lambda_0 = 2$ is a base rate (approximately the average number of
goals per team per 90-minute international match), $a_i$ is the attack
strength of team $i$ (equation~\ref{eq:attack}), and $d_i$ is the defence
strength of team $i$ (equation~\ref{eq:defence}).

The logic works as follows.  A strong attacking team ($a_H > 1$) inflates
$\lambda_H$.  A weak defensive team ($d_A > 1$) also inflates $\lambda_H$
because the stronger team's defence score is lower than 1 and the weaker
team's defence score is greater than 1.  The product neatly encodes
both effects simultaneously.

\textbf{Worked example.}  Suppose Argentina ($a \approx 1.38$,
$d \approx 0.62$) plays New Zealand ($a \approx 0.78$, $d \approx 1.22$):

\begin{align*}
  \lambda_{\text{ARG}} &= 2 \times 1.38 \times 1.22 \approx 3.37\\
  \lambda_{\text{NZL}} &= 2 \times 0.78 \times 0.62 \approx 0.97
\end{align*}

Argentina is expected to score around 3.4 goals, New Zealand under 1.

\subsection{Group-Stage Simulation}

In the group stage draws are allowed, so the code simply samples:

\begin{align*}
  g_H &\sim \Poi(\lambda_H)\\
  g_A &\sim \Poi(\lambda_A)
\end{align*}

and records the raw scoreline.

\begin{lstlisting}
def simulate_match_poisson(match, team_stats):
    home = team_stats[match.home]
    away = team_stats[match.away]
    lambda_home = expected_goals(home, away)
    lambda_away = expected_goals(away, home)
    home_goals = np.random.poisson(lambda_home)
    away_goals = np.random.poisson(lambda_away)
    return {'home_goals': home_goals, 'away_goals': away_goals}
\end{lstlisting}

\subsection{Knockout-Stage Simulation}

In the knockout stage a winner must emerge.  If the scores are level
after 90 minutes, extra time is simulated using a scaled-down Poisson
rate of $0.3\lambda$ (approximately 30 extra minutes relative to 90):

\begin{align}
  \delta_H &\sim \Poi(0.3 \, \lambda_H)\\
  \delta_A &\sim \Poi(0.3 \, \lambda_A)
\end{align}

If the match is still drawn after extra time, a penalty shoot-out is
resolved by a single Bernoulli trial.  The probability that the home
team wins the shoot-out is:

\begin{equation}
  p_{\text{pen}} = 0.5 + 0.1 \cdot
  \frac{r_H - r_A}{r_{\max} - r_{\min}}
  \label{eq:penalty}
\end{equation}

where $r_H$ and $r_A$ are the FIFA points of the home and away teams
respectively, and $r_{\max} - r_{\min} = 1877.27 - 1275.58$ is the
maximum observed spread across all teams in the tournament.  This gives a
slight advantage (up to $\pm 10\%$) to the higher-ranked team, but keeps
the result fundamentally uncertain.

\begin{lstlisting}
if home_goals == away_goals:
    extra_time = True
    extra_home = np.random.poisson(lambda_home * 0.3)
    extra_away = np.random.poisson(lambda_away * 0.3)
    home_goals += extra_home
    away_goals += extra_away

    if home_goals == away_goals:
        penalties = True
        ranking_benefit = 0.1*(home['Points'] - away['Points']) \
                          / (1877.27 - 1275.58)
        if random.random() < 0.5 + ranking_benefit:
            home_goals += 1
        else:
            away_goals += 1
\end{lstlisting}

% ────────────────────────────────────────────────────────────────────────────
\section{Step 3 --- Tournament Structure}
\label{sec:tournament}

\subsection{Groups}

The 2026 World Cup features 48 teams split into 12 groups of four (groups
A through L).  Within each group every team plays every other team once,
giving $\binom{4}{2} = 6$ matches per group and $12 \times 6 = 72$
group-stage matches in total.

\begin{lstlisting}
def generate_group_fixtures(groups):
    return {
        name: [Match(h, a) for h, a in combinations(teams, 2)]
        for name, teams in groups.items()
    }
\end{lstlisting}

\subsection{Group Standings}

Teams are ranked within each group by the standard FIFA tiebreak rules:

\begin{enumerate}[noitemsep]
  \item Points ($W = 3$, $D = 1$, $L = 0$)
  \item Goal difference ($\text{GF} - \text{GA}$)
  \item Goals scored ($\text{GF}$)
\end{enumerate}

The top two teams in each group qualify automatically.

\subsection{Third-Place Qualification}

The 2026 format advances the eight best third-placed teams to the
Round of 32.  The code collects all twelve third-placed teams, sorts them
by the same tiebreak criteria as above, and selects the top eight.

\subsection{The Knockout Bracket}

The Round of 32 consists of 16 matches following the official 2026 FIFA
bracket structure: group winners and runners-up are paired according to
the published draw, while the eight best third-placed teams fill the
remaining slots in ranked order.  Subsequent rounds are determined by
pairing consecutive match winners:

\[
  \text{Round}_{n+1} = \bigl\{\,
    \text{Match}(\text{winner}_{2i-1},\, \text{winner}_{2i})
    \mid i = 1, 2, \ldots
  \,\bigr\}
\]

The knockout rounds are: Round of 32 (16 matches), Round of 16
(8 matches), Quarter-finals (4 matches), Semi-finals (2 matches),
and the Final (1 match).

\begin{lstlisting}
def advance_round(matches):
    winners = [m.winner for m in matches]
    return [Match(winners[i], winners[i+1])
            for i in range(0, len(winners), 2)]
\end{lstlisting}

% ────────────────────────────────────────────────────────────────────────────
\section{Step 4 --- Monte Carlo Simulation}
\label{sec:montecarlo}

\subsection{The Idea}

Let $W_i$ be the event ``team $i$ wins the tournament''.  Its probability
$\Prob(W_i)$ is analytically intractable because it depends on the
random outcomes of up to seven matches, each of which determines the
opponents in subsequent rounds.  The Monte~Carlo estimator replaces this
with a frequency count:

\begin{equation}
  \hat{p}(W_i) = \frac{\text{number of simulations in which team } i
  \text{ wins}}{N}
  \label{eq:mc}
\end{equation}

By the Law of Large Numbers, $\hat{p}(W_i) \to \Prob(W_i)$ as
$N \to \infty$.

\subsection{Variance and Confidence}

For a true win probability $p$, the standard error of the Monte~Carlo
estimate is

\begin{equation}
  \sigma_{\hat{p}} = \sqrt{\frac{p(1-p)}{N}}.
\end{equation}

At $N = 100\,000$ simulations and $p = 0.2$ (a plausible win probability
for a strong team), this gives $\sigma_{\hat{p}} \approx 0.0013$, or
roughly $\pm 0.13\%$ --- more than sufficient for sweepstake purposes.

\subsection{What is Recorded}

For each run the simulator records:

\begin{itemize}[noitemsep]
  \item The champion.
  \item The finalist (runner-up).
  \item How far every team progressed (e.g.\ ``Quarter-finals'').
  \item The full scoreline of every group and knockout match, including
        whether the match went to extra time or penalties.
\end{itemize}

After $N$ runs, summary statistics are aggregated into several CSV files:
\texttt{summary.csv} (win percentages per team), \texttt{champions.csv}
(champion per run), \texttt{knockout\_results.csv},
\texttt{group\_results.csv}, \texttt{deep\_runs.csv}, and three
scoreline frequency files.

\begin{lstlisting}
N_SIMULATIONS = 100000

for run_id in range(1, N_SIMULATIONS + 1):
    group_fixtures = generate_group_fixtures(groups)
    simulate_group_stage_silent(group_fixtures)

    r32 = build_round_of_32(group_fixtures)
    ko_results = simulate_knockout_silent(r32)

    champion = ko_results['champion']
    champion_counts[champion] += 1
    ...
\end{lstlisting}

% ────────────────────────────────────────────────────────────────────────────
\section{Data Flow Summary}

The full pipeline can be summarised as follows:

\begin{enumerate}
  \item \textbf{Input.} \texttt{team\_rankings.csv} provides FIFA
        ranking points for each of the 48 teams.
  \item \textbf{Strength derivation.} \texttt{AttackDefenceStrength.py}
        converts each team's points into an attack strength $a_i$ and a
        defence strength $d_i$ via the linear map
        (equations~\ref{eq:attack}--\ref{eq:defence}), and writes the
        result to \texttt{team\_stats.csv}.
  \item \textbf{Match simulation.} For any pair of teams, \texttt{MatchSim.py}
        computes expected goals $\lambda_H$ and $\lambda_A$
        (equations~\ref{eq:lambdaH}--\ref{eq:lambdaA}), draws goals from
        $\Poi(\lambda_H)$ and $\Poi(\lambda_A)$, and resolves ties in
        knockout matches via scaled extra-time draws and a
        ranking-weighted Bernoulli trial (equation~\ref{eq:penalty}).
  \item \textbf{Tournament structure.} \texttt{Tournament.py} encodes the
        full 2026 bracket: 12 round-robin groups, best-eight-third-place
        qualification, and a 16-match knockout bracket through to the
        final.
  \item \textbf{Monte Carlo aggregation.} \texttt{TourSim.py} runs
        $N = 100\,000$ independent simulations, tallies results, and
        writes summary CSVs including win percentages and scoreline
        frequency tables.
\end{enumerate}

% ────────────────────────────────────────────────────────────────────────────
\section{Assumptions and Limitations}

\begin{description}[leftmargin=2em]
  \item[No home advantage.] The model labels teams ``home'' and ``away''
        for coding convenience but applies no home-ground boost.  In
        practice, the three host nations (USA, Canada, Mexico) might
        benefit from crowd support.
  \item[Static team strengths.] Strengths are computed once from the
        FIFA standings and held fixed throughout.  Injuries, squad
        rotation, and tournament form are not modelled.
  \item[Independent goals.] The Poisson model assumes the two teams' goal
        tallies are independent, which is a known simplification---teams
        often change tactics once they take the lead.
  \item[Linear FIFA-to-strength mapping.] The transformation
        (equation~\ref{eq:strength}) is a deliberate simplification.
        More sophisticated models (e.g.\ Dixon--Coles) estimate
        team-specific attack and defence parameters from historical match
        data rather than inferring them from a single rating.
  \item[Bracket simplification.] The official 2026 third-place slot
        allocation rules are more complex than the code's ranked-order
        assignment.  The code documents this explicitly as a simplification.
\end{description}

% ────────────────────────────────────────────────────────────────────────────
\section{Mathematical Summary}

For quick reference, the key equations are:

\begin{align}
  a_i &= 1 + \frac{r_i - 1500}{1000} \tag{attack strength}\\[4pt]
  d_i &= 2 - a_i \tag{defence strength}\\[4pt]
  \lambda_H &= 2 \cdot a_H \cdot d_A \tag{home expected goals}\\[4pt]
  \lambda_A &= 2 \cdot a_A \cdot d_H \tag{away expected goals}\\[4pt]
  g_H \mid \lambda_H &\sim \Poi(\lambda_H) \tag{home goals}\\[4pt]
  g_A \mid \lambda_A &\sim \Poi(\lambda_A) \tag{away goals}\\[4pt]
  p_{\text{pen}} &= 0.5 + 0.1 \cdot \frac{r_H - r_A}{601.69}
  \tag{penalty win probability}\\[4pt]
  \hat{p}(W_i) &= \frac{1}{N}\sum_{n=1}^{N}
    \mathbf{1}[\text{team } i \text{ wins run } n]
  \tag{Monte Carlo estimator}
\end{align}

% ────────────────────────────────────────────────────────────────────────────
\section*{Acknowledgements}

The predictor was written by Ben Jordan and Jonathan Glanfield to
accompany a sweepstake for the 2026 FIFA World Cup held jointly by
the United States, Canada, and Mexico.  The source code is publicly
available at
\url{https://github.com/JonathanG04157/worldcup}.

\end{document}